%% Dans ce fichier, les �ventuels accents sont cod�s
%% en ISO-8859-15 (aka latin9)
%% La ligne suivante donne trois exemples 
%% e accent aigu : �, e accent grave : �, e-dans-l'o : �
%%%%%ENTREES%%%%
\begin{tabarg}
\multicolumn{3}{l}{\titretab {Entr\'ees}}\\
\hline
$X$ & Vecteur, \texttt{float} & 1�re variable\\
$Y$ & Vecteur, \texttt{float} & 2�me variable\\
$F_e$ & Coeff., \texttt{float} & Fr�quence d'�chantillonnage de $X$ et $Y$ (\$Arg1.sample) ($Hz$)\\
$T_{int}$ & Coeff., \texttt{float} & Temps d'int�gration de la covariance  ($s$)\\
$Tlag_{min}$ & Coeff., \texttt{float} & D�calage temporel minimal � appliquer entre $X$ et $Y$ (s)\\
$Tlag_{max}$ & Coeff., \texttt{float} & D�calage temporel maximal � appliquer entre $X$ et $Y$ (s)\\
\end{tabarg}

\medskip
%%%%%SORTIES%%%%
\begin{tabarg}
\multicolumn{3}{l}{\titretab {Sorties}}\\
\hline
$Cov_{XY}$ & Vecteur, \texttt{float} & Covariance entre $X$ et $Y$ \\
$Corr_{XY}$ & Vecteur, \texttt{float} & Corr�lation entre  $X$ et $Y$ \\
$T_{valid}$ & Vecteur, \texttt{float} & Temps d'int�gration r�el de la covariance � partir 
du nombre de donn�es valides  ($s$)\\
$Tlag$ & Vecteur, \texttt{float} & D�calage temporel qui maxmimise la covariance  ($s$)\\
\end{tabarg}
